%! Periodic Phenomena — Unit 3, Lesson 3.1 — Student Packet Cover Sheet
\documentclass[10pt]{article}
\usepackage{precalculus-article}
\usepackage{precalculus-boxes}
\usepackage{ltablex}
\keepXColumns

\begin{document}

% Full-bleed navy banner
\begin{tikzpicture}[remember picture, overlay]
  \fill[navy] (current page.north west)
    rectangle ([yshift=-0.9in]current page.north east);
\end{tikzpicture}
\vspace{-0.8in}

\begin{center}
  {\color{white}\textbf{\LARGE Precalculus}}\\[4pt]
  {\color{skymid}\large\textbf{Unit 3: Trigonometric and Polar Functions}}\\[6pt]
  {\color{white}\textbf{\large Lesson 3.1 \quad Periodic Phenomena}}
\end{center}

\vspace{0.22in}
\namedateperiod
\vspace{0.18in}

\begin{learningtargetbox}
\small
\begin{enumerate}[leftmargin=1.5em, itemsep=5pt]
  \item I can identify a periodic relationship from a verbal or tabular description
    and estimate its period. \hfill\textit{(LO 3.1.A, 3.1.B)}
  \item I can construct a graph of a periodic relationship from one cycle of data.
    \hfill\textit{(LO 3.1.A)}
  \item I can describe key characteristics of a periodic function --- period,
    intervals of increase/decrease, and concavity --- and explain why each
    characteristic repeats in every period. \hfill\textit{(LO 3.1.B)}
\end{enumerate}
\end{learningtargetbox}

\vspace{0.14in}

\begin{tocbox}
\small
\renewcommand{\arraystretch}{1.5}
\begin{tabularx}{\linewidth}{@{} c l X r @{}}
\rowcolor{sky}
\textbf{\#} & \textbf{Component} & \textbf{Description} & \textbf{Score} \\
\midrule
1 & Warm-Up        & Spiral review: function notation, table patterns, rate of change & \blank{1.2cm} \\
2 & Guided Notes   & Periodic relationships, period, constructing graphs from one cycle & \blank{1.2cm} \\
3 & Group Activity & Identify and describe periodic patterns in context (tiered R/A/E) & \blank{1.2cm} \\
4 & Exit Ticket    & Identify period from a table; describe a repeating characteristic & \blank{1.2cm} \\
5 & Homework       & Mixed: identify, estimate period, AP-style MC & \blank{1.2cm} \\
\midrule
  & \multicolumn{2}{r}{\textbf{Total}} & \blank{1.2cm} \\
\end{tabularx}
\end{tocbox}

\end{document}
